\title{Adapting Language Models for Zero-shot Learning by Meta-tuning on Dataset and Prompt Collections}

\begin{abstract}
Large pre-trained language models (LMs) such as GPT-3 have acquired a surprising ability to perform zero-shot learning. For example, to classify sentiment without any training examples, we can ``prompt" the LM with the review and the label description ``\textit{Does the user like this movie?}", and ask whether the next word is ``\textit{Yes}" or ``\textit{No}".
However, the next word prediction training objective is still \textbf{misaligned} with the target zero-shot learning objective. To address this weakness, we propose meta-tuning, which directly optimizes the zero-shot learning objective by fine-tuning pre-trained language models on a collection of datasets. We focus on classification tasks, and construct the meta-dataset by aggregating 43 existing datasets and annotating 441 label descriptions in a question-answering (QA) format. When evaluated on unseen tasks, meta-tuned models outperform a same-sized QA model and the previous SOTA zero-shot learning system based on natural language inference. Additionally, increasing parameter count from 220M to 770M improves AUC-ROC scores by 6.3\%, and we forecast that even larger models would perform better. Therefore, measuring zero-shot learning performance on language models out-of-the-box might underestimate their true potential, and community-wide efforts on aggregating datasets and unifying their formats can help build models that answer prompts better.
\end{abstract}

\section{Introduction}

The goal of zero-shot classification (ZSC) is to classify textual inputs using label descriptions without any examples \cite{yin-etal-2019-benchmarking}. Large language models - whose only training objective is to predict the next word given the context - have acquired a surprising ability to perform ZSC \cite{radford2019language, brown2020language, le-scao-rush-2021-many}. For example, to classify whether the sentence ``\textit{This movie is amazing!}" is positive, we can prompt the language model with the context ``\textit{Review: This movie is amazing! Positive Review? \_\_\_ }", and check whether the next word is more likely to be ``\textit{Yes}" or ``\textit{No}" \cite{Zhao2021Calibrate}. To convert ZSC into a language modeling (LM) task that an LM model is likely to perform well, many recent works focus on finding better prompts \cite{shin-etal-2020-autoprompt, schick2020exploiting, schick2020s, gao-etal-2021-making}.

However, the LM training objective is correlated but still misaligned with the target objective to answer prompts.
\iffalse
For example, to classify review sentiment, we can ask the model to answer ``\textit{Yes/No}" to the question ``\textit{Does the user like this movie}?" (Figure \ref{fig:fig1} (a)). UnifiedQA \cite{khashabi-etal-2020-unifiedqa}, a model trained to answer generic questions and initialized with T5 (770 million parameters) \cite{raffel2019exploring}, can achieve zero-shot accuracy 92\% on SST-2 \cite{socher2013recursive} by answering this question, while GPT3 (175 B parameters) only obtains 80\% accuracy \cite{Zhao2021Calibrate}. Being adapted to QA makes a 200x smaller model answer prompts better. 
\fi
Our work addresses this weakness by directly optimizing the zero-shot classification objective through fine-tuning (Section \ref{sec:training}).
This requires us to 1) unify different classification tasks into the same format, and 2) gather a collection of classification datasets and label descriptions (prompts) for training (Section \ref{main-sec:data}).
Since we fine-tune our model on a meta-dataset, we name our approach meta-tuning.

We focus on binary classification tasks and unify them into a ``\textit{Yes}"/``\textit{No}" QA format \cite{clark-etal-2019-boolq, mccann2018natural}, where the input is provided as the context and the label information is provided in the question (Figure \ref{fig:fig1} (a)). Using this format, we gathered a diverse set of classification datasets from 43 different sources listed on Kaggle, SemEval, HuggingFace, and other papers. These tasks range from hate speech detection, question categorization, sentiment classification to stance classification, etc, and the genre ranges from textbooks, social media, to academic papers, etc. In total, these datasets contain 204 unique labels, and we manually annotated 441 label descriptions (Figure \ref{fig:data-collection}).

To evaluate ZSC, we need to define what counts as a task that the model has not seen during training time. While prior work considers different notions of ``unseen" by disallowing the same label or the same dataset to appear during training, our work defines ``unseen" more harshly by disallowing similar datasets. For example, we consider AG News topic classification dataset \cite{10.5555/2969239.2969312} and the topic classification dataset from \citet{yin-etal-2019-benchmarking} to be similar, even though their sources and label spaces are different.

Meta-tuning improves ZSC over UnifiedQA for most labels (Figure \ref{fig:fig1} (c)). Moreover, larger models are better, and hence we forecast that meta-tuning would work for even larger models. We also find that the performance can be slightly improved by training on datasets similar to the test dataset, ensembling different label descriptions, or initializing with a QA model  (Section \ref{sec:core-results}). 
All of our findings reliably hold under different robustness checks (Section \ref{sec:robust}), and our approach outperforms the previous SOTA \citet{yin-etal-2019-benchmarking} using the same pre-training method (Section \ref{sec:comparison}).

Our results suggest two promising future directions (Section \ref{sec:implications}).
First, large language models' (e.g. GPT-3) potential for zero-shot learning, as currently measured by context-prompting, might have been broadly underestimated; meta-tuning might significantly improve their performance. Second, community-wide efforts on aggregating and unifying datasets can scale up training and evaluation for zero-shot learning models. On the flip side, however, the meta-tuning approach might incentivize providers of LM inference APIs to collect prompts from users, hence potentially leading to security, privacy, and fairness concerns at a greater scale (Section \ref{sec:ethics}).

\paragraph{Contributions} To summarize, we 1) curate a dataset of classification datasets with expert annotated label descriptions. 2) demonstrate a simple approach to train models to perform zero-shot learning, and 3) identify several factors that improve performance; in particular, larger pretrained models are better.
\footnote{Code and data available here: \url{https://github.com/ruiqi-zhong/Meta-tuning}.}

\iffalse
We list several implications of our work in Section \ref{sec:implications} and \ref{sec:ethics}. 
First, since GPT-3 is not specialized in answering prompts, current measurements on its intelligence \cite{hendrycks2021aligning, hendrycks2021measuring} might significantly underestimate its true capability. Second, our work is still far from exhausting all NLP classification tasks, and the community might benefit from a shared effort of aggregating datasets and unifying their formats. At last, the meta-tuning approach might incentivize providers of language model inference APIs to collect prompts from the users, potentially leading to security, privacy and fairness concerns at a greater scale.
\fi

\iffalse To conclude, we contribute

\begin{itemize}
    \item a dataset of classification datasets with hand-written label descriptions.
    \item a meta-tuning approach that makes model specialized in zero-shot classification.
    \item a thorough discussion of the implication of our approach. 
\end{itemize}

\fi \iffalse Is it possible that the model only learns some shallow heuristics from meta-tuning, for example, that the word ``exciting" is more frequent in positive reviews? This is unlikely, since Section \ref{sec:shallw} finds that larger models are better for most label descriptions. Since we believe that smaller and larger models are equally capable of learning shallow heuristics, larger models must have used more complicated features during meta-tuning to perform better. Therefore, we forecast that meta-tuning would also be effective on even larger pre-trained models.

Another concern is that the UnifiedQA is initialized with T5 \cite{raffel2019exploring}, which is trained on the open web and might have seen the datasets we gathered; therefore, our observation could be an artifact that larger models are better at memorizing the datasets in the pretraining corpus \cite{carlini2020extracting, sagawa2020investigation}. To investigate this possibility, Section \ref{sec:memorization} compares different sizes of BERT \cite{devlin-etal-2019-bert}, which is trained on Wikipedia and Book Corpus \cite{zhu2015aligning}. We still find that larger models are better and most AUC-ROC scores are well above 0.5. Hence, our conclusion is robust and the model performs well not only because of memorization.
\fi

\section{Data} \label{main-sec:data}
We gather a wide range of classification datasets and unify them into the ``\textit{Yes}"/``\textit{No}" question answering format for binary classification. Then we group similar datasets together to determine what counts as unseen tasks during evaluation. 

\paragraph{Gathering classification datasets} We collect classification datasets from
Kaggle\footnote{\url{https://www.kaggle.com}}, 
Huggingface \cite{wolf-etal-2020-transformers},
SemEval\footnote{\url{https://semeval.github.io} }, and other papers. We looked through these sources and only considered English classification datasets. We also skipped the tasks that we felt were already better represented by other datasets in our collection. Then we manually examined a few examples in each remaining dataset to make sure it seemed plausibly clean.

The goals of these classification datasets include, but are not limited to sentiment classification (IMDB Reviews, \citet{maas-EtAl:2011:ACL-HLT2011}), topic classification (AG News, \citet{10.5555/2969239.2969312}), grammaticality judgement (CoLA, \citet{warstadt2018neural}), paraphrase detection (QQP\footnote{\url{https://www.kaggle.com/c/quora-question-pairs}}), definition detection (SemEval 2020 Task 6, \citet{spala-etal-2019-deft}), stance classification (SemEval 2016 Task 6, \citet{mohammad-etal-2016-semeval}),  etc. The genre includes academic papers, reviews, tweets, posts, messages, articles, and textbooks. The comprehensive list of datasets is in Appendix \ref{appendix-sec:dataset}.
Overall, we aim for a high diversity of tasks and genres by building upon what the broader research community has studied. Our approach is complementary to that of \citet{weller-etal-2020-learning}, which asks turkers to generate tasks, and that of \citet{mishra2021natural}, which generates tasks by decomposing existing templates used to construct reading comprehension datasets. The concurrent work of \citet{bragg2021flex} unifies the evaluation for few-shot learning; their zero-shot evaluation setup is the closest to ours, and they used templates and verbalizers \cite{schick2020exploiting} to specify the semantics of a task. 

Some of our datasets are noisy and not peer reviewed, or contain tasks that are too complicated (e.g. Multi-NLI, \citet{williams-etal-2018-broad}) for ZSC. To make our evaluation more informative, we only include them for training but not testing. We make these decisions before running our experiments in Section \ref{sec:results} to prevent selection bias. 

\paragraph{Unifying the dataset format}

We convert each classification dataset into a \textit{``Yes"/``No"} question answering format and provide label information in the question. For each label, we annotate 1-3 questions. If the label is null (for example, a text that does not express a particular emotion in an emotion classification dataset), we skip this label. Three of the authors\footnote{One of them is a graduate student and the other two are undergrads; all of them study Computer Science and have taken an NLP class.} manually annotated 441 questions for 204 unique labels, and each question is proofread by at least another author. See Figure \ref{fig:data-collection} for a concrete example, and Figure \ref{fig:example-descriptions} for some representative label descriptions. 

Additionally, some datasets contain thousands of labels \cite{chalkidis-etal-2019-large, allaway-mckeown-2020-zero}. In this case, we use templates to automatically synthesize label descriptions and exclude them from evaluation.

\paragraph{Grouping similar datasets} Our goal is to test the models' ability to generalize to tasks that are different enough from the training tasks. Therefore, at test time, we need to exclude not only the same dataset that appeared in the meta-tuning phase, but also ones that are similar.

This poses a challenge: whether two datasets perform the same task involves subjective opinion, and there is no universally agreed definition. On one extreme, most datasets can be counted as dis-similar tasks, since they have different label spaces and input distributions. On the other extreme, all datasets can be considered the same task, since they can all be unified into the question answering format.

To tackle this challenge, we create a set of tags, each describing a dataset property. The set of tags includes \textit{domain classification}, \textit{article}, \textit{emotion}, \textit{social-media}, etc, and the full set of them can be seen in Appendix \ref{appendix-sec:dataset-property}.
Then we define the two datasets to be similar if they are associated with the same set of tags, and prohibit the model to learn from one and test on the other. For example, our work considers the topic classification datasets from \citet{10.5555/2969239.2969312} (AG News) and \citet{yin-etal-2019-benchmarking} to be similar since they both classify topics for articles, even though their sources and label spaces are different. Some example dataset groups can be seen in Figure \ref{fig:task-group}.

Nevertheless, our procedure is not bullet-proof and one can argue that our notion of unseen tasks, though harsher than prior works \cite{yin-etal-2019-benchmarking, pushp2017train}, is still lenient. Therefore, as additional robustness checks, for each dataset we evaluate, we manually identify and list the most relevant dataset that is allowed during training in Appendix \ref{appendix-sec:most-relevant} . For example, the most relevant dataset to the IMDB review sentiment classification dataset is the emotion classification dataset from \citet{yin-etal-2019-benchmarking}, which classifies the input text into 9 emotions, such as ``\textit{joy}", ``\textit{surprise}", ``\textit{guilt}", etc. We consider the emotion classification dataset to be relevant, since sentiment classification often involves identifying emotions. However, one can also argue that they are different tasks: their input and label spaces are different, and sadness can be caused by a great tragedy, or a bad movie that wastes the users' time. The comprehensive list of label descriptions grouped by dataset similarity is in Appendix \ref{appendix-sec:label-descriptions}.

In total, we spend around 200 hours to collect this dataset. This time estimate includes skimming through the dataset repos and recent NLP papers, writing programs to download the datasets and unify their format, annotating label descriptions, performing quality controls, and documenting the collection process.

\section{Metrics} \label{sec:metrics}
To reliably aggregate performance across different datasets and present as much information as possible, we report a set of descriptive statistics and provide visualizations whenever we compare two models. We generally do not reduce a model's performances on different datasets into one scalar quantity and compare this number only.

\paragraph{Descriptive statistics} For each label description (question), we calculate the AUC-ROC score
\footnote{We do not evaluate F-score or accuracy, since they are very sensitive to the decision cutoff, and usually additional calibration is needed \cite{Zhao2021Calibrate}.} by treating the ``\textit{Yes}" answer as the positive class. After calculating the AUC-ROC score for each label, we calculate the following set of descriptive statistics to compare two models. Suppose that model $Y$ is hypothetically better than $X$. Denoting $\Delta$ as the change of AUC-ROC of a label description from $X$ to $Y$, we can summarize how $\Delta$ is distributed across the set of label descriptions with the following statistics:

\begin{itemize}
    \item $\mathbb{E}[\Delta]$: the average change in AUC-ROC.
    \item $\mathbb{P}[\Delta > t]$: the fraction of label descriptions where the change is over the threshold $t$. 
    \item $\mathbb{P}[\Delta < -t]$: the fraction of label descriptions where the change is less than $-t$. 
    \item $Std[\Delta]$: the standard deviation of the change.
\end{itemize}

In the main paper, we weight each label description equally in this distribution to calculate the above statistics. We may also weight each label or dataset equally, and the corresponding results are in Appendix \ref{appendix-sec:comprehensive-results}.
To make sure our conclusions are robust, we consider one model to be better only when $\mathbb{E}[\Delta] > 0$ and $\mathbb{P}[\Delta > t] > \mathbb{P}[\Delta < -t]$ for all $t \in \{1\%, 5\%, 10\%\}$, under all three types of weighting. In other words,  we claim that one model is better than the other only when 12 conditions simultaneously hold. 

\paragraph{Visualizations} We use scatter plots to visualize and compare the performance of two models, where each dot represents a label description, its x-value represents the AUC-ROC score of the model $X$, and its y-value represents that of $Y$. If most dots are above the identity line $y=x$, the model $Y$ is better than $X$. 

The descriptive statistics and the visualizations are explained in Figure \ref{fig:metrics-explained}.

\section{Model} \label{sec:training}

\paragraph{Architecture} 
We format the inputs to the model in the same way as UnifiedQA \cite{khashabi-etal-2020-unifiedqa}, which concatenates the context to the question and adds a ``\textit{[SEP]}" token in between. Then we feed the concatenated input into the T5 encoder and produce the answer score by normalizing the ``\textit{Yes}"/``\textit{No}" probability of the first decoded token. Unless otherwise noted, we initialize our model with T5-Large (770 Million parameters). We sometimes compare to or initialize with the UnifiedQA model \cite{khashabi-etal-2020-unifiedqa}, which is trained on a wide range of question answering datasets. For a fair comparison, we use the UnifiedQA model initialized with T5-Large as well. To meta-tune non-Seq2Seq pre-trained models, such as BERT \cite{devlin-etal-2019-bert} or RoBERTa \cite{liu-etal-2019-robust}, we add an MLP layer on top of the pooled output/``\textit{[CLS]}" token to classify between ``\textit{Yes}"/``\textit{No}". 
We leave the improvement on model architectures \cite{ye2021zero, li2021prefix, lester2021power} and training objectives \cite{murtydreca, Yin2020Meta-Learning} for future work. 

\paragraph{Meta-tuning}
We create a training distribution that balances between datasets, label descriptions, and ``\textit{Yes}"/``\textit{No}" answers. To create the next training datapoint for meta-tuning, we select a dataset from the training split uniformly at random (u.a.r.); then we select a label description (question) u.a.r. and with 50\% probability select a textual input with the answer ``\textit{Yes}"/``\textit{No}".
To prevent over-fitting, we do not train on any combination of label description and textual input twice. Unless otherwise noted, we meta-tune the model for 5000 steps and use batch size 32. We did not tune any hyper-parameters or training configurations since they work well during our first attempt. To evaluate ZSC performance on each dataset, we leave out one group of similar datasets as the evaluation set and train on the rest. Altogether, the experiments take around 250 GPU hours on Quadro 8000. 

\section{Results} \label{sec:results}

\subsection{Hypotheses and Conclusions} \label{sec:core-results}

We investigate and validate the following hypotheses, sorted by importance in descending order.

\begin{itemize}
    \item Meta-tuned models outperform general question answering models in zero-shot classification.
    \item Larger pre-trained models are better.
    \item Pre-training does the heavy lifting.
    \item Performance can be improved by training on similar datasets, initializing with a QA model, or ensembling label descriptions.
    \item Early stopping is crucial to performance.
\end{itemize}

\paragraph{Meta-tuned models are better.} We compare a meta-tuned T5-Large model (770 M parameters)\footnote{This model is initialized with T5, not UnifiedQA.} with the same-sized UnifiedQA model \cite{khashabi-etal-2020-unifiedqa} out of the box. Relevant descriptive statistics can be seen in the first row of Table \ref{tab:main-tab} and Figure \ref{fig:main_large} (a). Adapting the model for ZSC improves the average AUC-ROC by 3.3\%. 
\paragraph{Larger pre-trained models are better.} We compare T5-Base (220 Million parameters) against T5-Large (770 M). The statistics can be seen in the second row of Table \ref{tab:main-tab} and Figure \ref{fig:main_large} (b). Increasing the model size from 220 M to 770M improves the average AUC-ROC by 6.3\%. 

\paragraph{Pre-training does the heavy lifting.} In Figure (c) and the third row of Table \ref{tab:main-tab}, we compare pre-trained and random initializations, where the latter cannot beat the random baseline (average AUC-ROC 0.503). Hence, meta-tuning alone is far from enabling the model to perform ZSC. An intuitive interpretation is that the model already ``knows" how to perform ZSC after pre-training under the LM objective, and learns how to use this knowledge during meta-tuning.

\paragraph{Training on similar datasets improves performance.} Unlike before, we no longer avoid training on similar datasets from the same group. Instead, we perform straightforward leave-one-out cross-validation. The statistics can be seen in the fourth row of Table \ref{tab:main-tab} and Figure \ref{fig:main_large} (d), and it improves the average AUC-ROC by 0.7\%.
The performance gain is not as significant as increasing the model size or adapting for ZSC. We conjecture that it is because we have not collected enough datasets; otherwise, there might be more similar datasets, hence improving ZSC performance.

\paragraph{Ensembling label descriptions improves performance.} Instead of asking the model a single question for each label and obtain the probability of the answer being ``\textit{Yes}", we can average the probability obtained by asking multiple questions with the same meaning. This approach is different from traditional ensembling, which typically needs to store/train multiple models to average across them. The fifth row of Table \ref{tab:main-tab} and Figure \ref{fig:main_large} (e) verifies that ensembling descriptions improves performance slightly (0.7\% AUC-ROC score).

\paragraph{Initializing with UnifiedQA improves performance.} Figure \ref{fig:main_large} (f) and the sixth row of Table \ref{tab:main-tab} compare the UnifiedQA against against the T5 initialization. Initializing with UnifiedQA improves average AUC-ROC by 1.1\%.

\paragraph{Early stopping is crucial to performance.} If we train the model for too long, the model might simply ``memorize" that certain label descriptions correspond to certain training tasks, and the performance on unseen tasks may drop. To explore this possibility, we meta-tune our models for 100K steps, which is 20 times as long as our default setting and encourages the model to memorize the training tasks. We then evaluate them on the three benchmark zero-shot classification datasets by \citet{yin-etal-2019-benchmarking} (which we describe in more details in the next section). We calculate the average AUC-ROC across all label descriptions for each of the 3 datasets, and plot them in Figure \ref{fig:longrunning}.

The performance decreases \footnote{Kendall rank correlation coefficients are negative with $p < 0.005$ for topic and situation classification} as training continues. On the other hand, however, the performance drop of 3\% in AUC-ROC is not fatal and the model's performance is still much better than random guesses.

\subsection{Robustness Checks} \label{sec:robust}
We examine a series of additional results to make sure our conclusions are robust. The observed improvements in Table \ref{tab:main-tab} and Figure \ref{fig:main_large} might be caused by the improvement of a small number of labels that are annotated with more descriptions, or by the improvement on a dataset with more distinct labels. Appendix \ref{appendix-sec:other-uniform} compares the performance by assigning equal weights to each label/datasets.

To provide additional supporting evidence for our forecast that larger models are better, Appendix \ref{appendix-sec:t5-larger-better} compares a 60M-parameter model against a 220M-parameter model, and finds that the latter is much better. One concern, however, is that our models are initialized with T5 \cite{raffel2019exploring}, which is trained on the open web and might have seen the datasets we gathered. Therefore, larger models might be better simply because they are better at memorization \cite{sagawa2020investigation}. Appendix \ref{appendix-sec:bert-larger-better} addresses this by showing that larger models are also better with BERT initialization \cite{devlin-etal-2019-bert}, which is trained on Wikipedia and Book Corpus \cite{zhu2015aligning}.

We also report the models' performance on each dataset for readers' reference in Appendix \ref{appendix-sec:breakdown}.

\subsection{Comparison with \citet{yin-etal-2019-benchmarking}} \label{sec:comparison}

This section shows that our approach has higher performance than the zero-shot classification system built by \citet{yin-etal-2019-benchmarking}.
Their system ensembles several natural language inference models based on RoBERTA-Large (355M parameters, \citet{liu2020roberta}), and another model trained to categorize Wikipedia articles. It was evaluated on three classification datasets:

\begin{itemize}
    \item topic (10-way): classifies article domains, such as \textit{family \& relationship, education, sports}, etc. The metric is accuracy. 
    \item emotion (10-way): classifies emotion types, such as \textit{joy, anger, guilt, shame}, etc. The metric is label-weighted F1.
    \item situation (12-way): classifies disaster situations, e.g. \textit{regime change, crime \& violence}, and the resource they need, e.g. \textit{search \& rescue}. The metric is label-weighted F1.
\end{itemize}

We use the exact same evaluation metrics as in \citet{yin-etal-2019-benchmarking}, and the same label resolution strategy when the model answers ``Yes"\footnote{or ``Entailment" for natural language inference models.} for multi-label classification. Concretely, when the model predicts ``Yes" on multiple labels, the one with the highest probability is selected. For a fair comparison, we meta-tune RoBERTa of the same size and compare it with the highest performing model in \citet{yin-etal-2019-benchmarking} for each of the three datasets.

The results are in Table \ref{tab:benchmarking}, and our model has higher performance across all 3 datasets using the same pre-training method.

\section{Discussion and Future Directions} \label{sec:implications}

\paragraph{Main takeaways} We construct a dataset of classification datasets to adapt the language model for zero-shot classification via meta-tuning. The adapted model outperforms a general-purpose question answering model and the prior state of the art based on natural language inference. We forecast that meta-tuning would be more effective on larger models, and the current engineering ceiling for zero-shot learning might have been broadly under-estimated.

\paragraph{Aggregating and unifying datasets} The main bottleneck of our research is to manually gather a wide range of datasets and unify their format. The difficulties are: 1) we need to brainstorm and review the NLP literature extensively to decide what new tasks to look for; 2) different datasets encode their data in different formats, and we need to write programs manually for each of them to convert to the desired format; 3) it is hard to tell the quality of a dataset purely by its provenance, and sometimes we need to examine the dataset manually. If we as a community can aggregate and unify datasets better, we could potentially train and evaluate zero-shot learning models at a larger scale.

\paragraph{Meta-tuning as a probe} There is a growing interest in measuring the intelligence \cite{hendrycks2021aligning, hendrycks2021measuring} or the few-shot learning ability \cite{brown2020language} of large language models like GPT-3. However, since these models are not adapted to answer those prompts \cite{DBLP:journals/corr/abs-2104-08315}, we suspect that its knowledge and true potential to perform few-shot learning is much higher than reported. Since pre-training does the heavy lifting and meta-tuning is unlikely to provide additional ZSC ability to the model, we can potentially first use meta-tuning as a probe to make them adapted to answering prompts before measuring their performance.

Still, to make this methodology rigorous, interpreting and controlling the strength of the probes will be an important future direction \cite{hewitt-liang-2019-designing}. For example, if the training set contains a prompt that is too similar to the prompt to be tested, the probe will be meaningless.

\paragraph{Beyond Shallow Correlations} One possibility is that the model only learns shallow statistical correlations from meta-tuning rather than ``more sophisticated reasoning skills". For example, the word ``exciting" might occur in positive reviews more. This is unlikely, given that larger models are consistently better than smaller or randomly initialized ones. To explain this performance gap, larger models must have learned to use more complicated features during meta-tuning.

\paragraph{Relation to Meta/Multitask-Learning} Our method is closely related to, but different from meta-learning \cite{yin2020meta, murtydreca} and multi-task learning \cite{DBLP:journals/corr/abs-2104-08835, aghajanyan2021muppet}. Both meta-learning and multitask-learning typically involve at least a few examples from the target task; in our setup, however, the model does not learn from any target task examples. The “meta” in our name does not mean “meta-learning”, but reflects the fact that our model learns from a meta-dataset of tasks. 

Nevertheless, our framework can be easily adapted to a few-shot learning setup, which enables the language model to learn to learn from in-context examples (see below). Since this approach models the learning process as a sequence classification problem, it can be seen as a form of meta-learning similar to \cite{ravi2016optimization}.

\paragraph{Annotating Prompts} Three of our authors annotated the label descriptions. Since they are all Computer Science major students who understand machine learning and natural language processing, they might not be representative of the final user population of this ZSC application. Annotating prompts that match the target user distribution will be an important research direction. 

Additionally, shorter and more natural descriptions sometimes fail to capture the exact semantics of the label. For example, in \citet{yin-etal-2019-benchmarking}, the description of the label ``medical" is ``\textit{people need medical assistance}"; or alternatively, it can be longer but more accurate: ``\textit{people need an allied health professional who supports the work of physicians and other health professionals}". How to scalably generate more accurate and detailed label descriptions without expert efforts will be another future direction. 

\paragraph{Optimizing Prompts} Our work is complementary to recent works that optimize the prompts to achieve better accuracy. Even if our meta-tuned model is specialized in answering prompts, it might still react very differently towards different prompts. For example, in the stance classification dataset \cite{barbieri-etal-2020-tweeteval}, we annotated two label descriptions (prompts) for the same label: ``Does this post support atheism?" and ``Is the post against having religious beliefs?". They have similar meanings, but the former has much lower accuracy than the later. We conjecture that this is because the model cannot ground abstract concepts like ``atheism".

\paragraph{Other extensions} We conjecture that meta-tuning can be extended to more diverse tasks beyond zero-shot binary classification. To extend to multi-label classification, we need to develop a procedure to resolve the labels when the model predicts positive for more than one labels. To extend to few-shot learning, we need to increase the context length to fit several training examples into the input, which requires a larger context window and hence more computational resources.

To extend to other sequence generation tasks, we need to collect a wide range of diverse sequence generation tasks to meta-tune the model, such as machine translation, summarization, free-form question answering, grammar correction, etc.

\end{document}
